\section{結果：愛媛（愛南町）}
\label{sec:ehime}

3 つ目の領域は、愛媛県南端の愛南町から上げる場合を想定したものである。
\texttt{sites.json} の「愛媛・高知」には愛南町と周辺の 8 地点が登録されており、
代表地点はそのうち\textbf{南レク御荘公園}（\ang{32.9642}N, \ang{132.5521}E）とした。
和歌山と違って打ち上げ地点が決まっているので、
北海道と同じく地点ごとの表（表\ref{tab:ainan-offshore}）も作れる。

対象は \ang{31.5}--\ang{34.5}N / \ang{132.0}--\ang{136.0}E、
\ang{0.375} 格子で 108 点。豊後水道の東半分から四国の南半分、土佐湾、
紀伊水道の西側までを含む。

\begin{whybox}{領域が打ち上げ地点を中心にしていない理由}
愛南町は、取得済みのデータ範囲（マスター箱、\ang{132}--\ang{150}E）の西端から
\ang{0.55}（約 \SI{51}{\kilo\metre}）しか離れていない。
西へ広げるにはデータの取り直しが要るので、
領域は地点を中心にせず、\textbf{東へ伸ばした}。
この地点から上げた機体は 1 年のほとんどで東へ流されるので、
統計が要るのも東側である。

代償が 2 つある。
1 つは、領域平均が地点より東（室戸岬や紀伊水道寄り）の風を多く含むことである。
ただし後で見るとおり、この領域の中では風がほとんど変わらないので（図\ref{fig:f6-ehime}）、
平均が地点の値から大きく離れることはない。
もう 1 つは、夏に西へ流れる機体がデータの外へ出やすいことで、
第 \ref{sec:fidelity} 章の軌道計算ではこれを数えて除外している
（\ref{sec:fidelity-error} 節）。
\end{whybox}

最寄り格子（\ang{32.78}N, \ang{132.38}E）までは \SI{26.5}{\kilo\metre} ある。
格子の対角のほぼ真ん中に地点が入るためで、\ang{0.375} 格子で起こりうる最大に近い。
それでも下で見るように領域内の風の差は小さいので、この距離が結論を変えることはない。

\subsection{年間を通した姿}

\begin{figure}[H]
\centering
\includegraphics[width=0.95\linewidth]{ehime/F0_sample_counts.pdf}
\caption[標本数]{各ビンに入った解析値の数。前半は必ず 15 日ぶんで 600 個、後半は月の長さで変わり 1 月は 640 個、2 月は閏年が 3 回あるので 532 個になる。この図が合わなければ取得か集計が壊れているので、他のどの図よりも先に確認する。}
\label{fig:f0-ehime}
\end{figure}

\begin{figure}[H]
\centering
\includegraphics[width=1.0\linewidth]{ehime/F1_alt_half_speed.pdf}
\caption[時期×高度の風速と風向]{上：平均風速。下：平均風向（吹いてくる向き）。風向には円環用のカラーマップを使っている。順次カラーマップを使うと \ang{0} と \ang{360} の境目に実在しない断層が現れるためである。灰色は定常度が 0.3 未満で塗りを外した領域で、そのビンではそもそも風向が定まっていない。}
\label{fig:f1-ehime}
\end{figure}

\begin{figure}[H]
\centering
\includegraphics[width=1.0\linewidth]{ehime/F3_profiles_by_month.pdf}
\caption[月別の高度プロファイル]{月ごとの風速プロファイル。線は各ビンの中央値、灰色の帯はその月の 2 ビンを通した p10--p90 の広がり。帯の幅が平均からのずれの大きさを表しており、線だけを見て計画を立ててはいけない。}
\label{fig:f3-ehime}
\end{figure}


愛媛の上空は、\textbf{和歌山とほぼ同じ気候}である。
\SI{200}{\hecto\pascal}（約 \SI{12.2}{\kilo\metre}）の平均風速は
冬 \SI{75.4}{\metre\per\second}（定常度 0.98）、春 \SI{49.6}{\metre\per\second}、
夏 \SI{21.9}{\metre\per\second}（定常度 0.55）、秋 \SI{39.4}{\metre\per\second}。
和歌山の冬 \SI{75.5}{\metre\per\second}、夏 \SI{22.7}{\metre\per\second} とほとんど変わらない。
いちばん強いのは 2 月前半の \SI{78.8}{\metre\per\second} である。

2 つの領域は緯度がほぼ同じ（\ang{33}--\ang{34}N 付近）で、
冬の亜熱帯ジェットの核がどちらの真上も通る。
経度が \ang{2} ほど違っても、ジェットの強さはほとんど変わらない。
北海道（冬 \SI{40.2}{\metre\per\second}）との差は、経度ではなく緯度の差である。

\subsection{成層圏}

\begin{figure}[H]
\centering
\includegraphics[width=1.0\linewidth]{ehime/F2_alt_half_constancy.pdf}
\caption[定常度]{定常度（ベクトル平均風速 ÷ スカラー平均風速）。1 に近いほど風向が安定し、0 に近いほどビンのなかで風向が入れ替わっている。淡い帯が成層圏の風の反転している高度で、季節とともに上下する。}
\label{fig:f2-ehime}
\end{figure}

\begin{figure}[H]
\centering
\includegraphics[width=1.0\linewidth]{ehime/F7-1_wind_roses.pdf}
\caption[風配図（前半）]{代表格子点の風配図、\textbf{月の前半}（1--15 日）。
平均と楕円では見せられない「分布そのもの」を示す。
ローブが 2 つに割れている場面が、平均風がほぼゼロで分散だけが大きくなる時期にあたる。}
\label{fig:f7-ehime-h1}
\end{figure}

\begin{figure}[H]
\centering
\includegraphics[width=1.0\linewidth]{ehime/F7-2_wind_roses.pdf}
\caption[風配図（後半）]{同じ格子点の\textbf{月の後半}（16 日--月末）。
前半（図\ref{fig:f7-ehime-h1}）と並べると、ほとんどの月で\textbf{2 枚はよく似ている}。
唯一 10 月だけ \SI{30}{\hecto\pascal} の最頻セクターが前半 \ang{90}、
後半 \ang{248} と入れ替わるが、どちらも頻度は 2 割前後しかない。
つまり向きが変わったというより、\textbf{その時期はどちらの半月も
風向が定まっていない}のであって、最頻セクターを読む意味が薄い。}
\label{fig:f7-ehime-h2}
\end{figure}

\SI{10}{\hecto\pascal}（約 \SI{31}{\kilo\metre}）は、
夏に \SI{18.5}{\metre\per\second}・定常度 0.99 の\textbf{東風}である。
冬と春の定常度は 0.38、0.36 しかなく、
最も低いのは 3 月後半（0.05）、1 月後半（0.13）、11 月前半（0.13）、
2 月前半（0.17）、3 月前半（0.22）である。
順位の入れ替わりはあるが、和歌山と同じ時期が並んでいる。
冬の \SI{10}{\hecto\pascal} が東風（1 月の領域平均 $u = \SI{-4.2}{\metre\per\second}$）
なのも北海道・和歌山と同じで、アリューシャン高気圧の西縁にあたる。

風配図（図\ref{fig:f7-ehime-h1}、\ref{fig:f7-ehime-h2}）の格子点は、
愛南町ではなく領域の中央（\ang{33.15}N, \ang{134.25}E、室戸岬付近）である。
この図は領域の中央で描く決まりにしてあり、3 領域とも同じである。

\subsection{領域の中で答えが変わるか}

\begin{figure}[H]
\centering
\includegraphics[width=1.0\linewidth]{ehime/F6_spatial_spread.pdf}
\caption[領域内のばらつき]{領域内での平均風速の幅（最大 $-$ 最小）。これが小さい高度・時期では、領域のどこで上げても風はほぼ同じで、代表 1 地点のプロファイルで語ってよい。大きいところではそうではない。}
\label{fig:f6-ehime}
\end{figure}


領域内での平均風速の幅は最大 \SI{13.3}{\metre\per\second} で、
和歌山（\SI{12.0}{\metre\per\second}）とほぼ同じ、
北海道（\SI{31.7}{\metre\per\second}）の半分以下である。
和歌山よりわずかに大きいのは、東西に \ang{4} と和歌山より \ang{1} 広いためである。

したがって\textbf{愛媛の領域も、代表 1 地点でおおむね語ってよい}。
領域平均が東へ偏っていること、最寄り格子まで \SI{26}{\kilo\metre} あることは、
この小ささのおかげで問題にならない。
実際、代表地点の最寄り格子で見ても \SI{200}{\hecto\pascal} は
冬 \SI{75.7}{\metre\per\second}、夏 \SI{21.4}{\metre\per\second} で、領域平均と 1 割も違わない。

\subsection{地図で見る}

\begin{figure}[H]
\centering
\includegraphics[width=1.0\linewidth]{ehime/M1-200_month_grid_200hPa.pdf}
\caption[200 hPa の月別分布図]{\SI{200}{\hecto\pascal}（約 \SI{12}{\kilo\metre}）の平均風を月ごとに並べたもの。塗りが風速、矢印が風向。ジェットが季節とともに動く様子が地図の上で見える。}
\label{fig:m1200-ehime}
\end{figure}

\begin{figure}[H]
\centering
\includegraphics[width=1.0\linewidth]{ehime/M1-030_month_grid_30hPa.pdf}
\caption[30 hPa の月別分布図]{同じく \SI{30}{\hecto\pascal}（約 \SI{24}{\kilo\metre}）。対流圏とは逆の季節変化をしており、夏に東風が強まる。}
\label{fig:m1030-ehime}
\end{figure}

\begin{figure}[H]
\centering
\includegraphics[width=1.0\linewidth]{ehime/M3-200_direction_map_200hPa.pdf}
\caption[風向の分布図]{\SI{200}{\hecto\pascal} の平均風向。円環用のカラーマップを使い、定常度 0.3 未満は灰色にしている。灰色の領域では平均風向に意味がない。}
\label{fig:m3200-ehime}
\end{figure}

\begin{figure}[H]
\centering
\includegraphics[width=1.0\linewidth]{ehime/M2_constancy_map.pdf}
\caption[定常度の分布図]{高度と季節ごとの定常度。濃いところでは風向が安定していて計画が立てやすく、淡いところではビンのなかで風向が入れ替わる。}
\label{fig:m2-ehime}
\end{figure}

\begin{figure}[H]
\centering
\includegraphics[width=0.92\linewidth]{ehime/M4-01_wind_ellipse_map.pdf}
\caption[風のばらつきを地図に]{1 月前半、\SI{200}{\hecto\pascal} の平均風ベクトルと、その先端に描いた $1\sigma$ の散らばり楕円。丸い楕円はどの向きにもばらつくこと、細長い楕円は主に強さがばらつくことを意味する。}
\label{fig:m401-ehime}
\end{figure}


\subsection{落下位置}

\begin{figure}[H]
\centering
\includegraphics[width=1.0\linewidth]{ehime/M5_drift_month_grid.pdf}
\caption[月別の平均落下位置]{各格子点から上げた場合の平均落下位置。\textbf{矢印は実寸}で、矢先がそのまま予測落下点である。色が平均の飛行距離。全パネルが同じ枠と同じ色スケールなので、月どうしを直接比べられる。}
\label{fig:m5-ehime}
\end{figure}

\begin{figure}[H]
\centering
\includegraphics[width=0.92\linewidth]{ehime/M7-01_drift_map.pdf}
\caption[1 月の落下位置]{1 月前半の平均落下位置を拡大したもの。枠は集計対象の領域で、矢先はその外へ出ていることが多い。}
\label{fig:m701-ehime}
\end{figure}

\begin{figure}[H]
\centering
\includegraphics[width=0.92\linewidth]{ehime/M7-07_drift_map.pdf}
\caption[7 月の落下位置]{同じく 7 月前半。1 月と比べると距離も向きも大きく違う。}
\label{fig:m707-ehime}
\end{figure}

\begin{figure}[H]
\centering
\includegraphics[width=1.0\linewidth]{ehime/M6_drift_percentile_map.pdf}
\caption[荒れた日の飛行距離]{飛行距離の 90 パーセンタイル。平均 \SI{150}{\kilo\metre} で p90 が \SI{400}{\kilo\metre} という状況と、平均 \SI{150}{\kilo\metre} だけを見て計画することは、運用上まったく違う。}
\label{fig:m6-ehime}
\end{figure}

\begin{figure}[H]
\centering
\includegraphics[width=0.92\linewidth]{ehime/M8_landing_ellipse_map.pdf}
\caption[代表地点からの落下分布]{1 地点から上げた場合の落下位置を季節別に重ねたもの。楕円は 50 / 90 / 95 \% の範囲で、各解析時刻ごとに求めた着地点の共分散から描いている。散らばりが 2 変量正規に従うという仮定は残る。灰色が陸、白が海である。}
\label{fig:m8-ehime}
\end{figure}


飛行時間は約 \num{137} 分、領域平均の飛行距離は年間平均 \SI{156}{\kilo\metre}。
和歌山（\SI{158}{\kilo\metre}）とほぼ同じである。

\begin{table}[H]
\centering
\caption{愛媛：飛行距離が最短・最長のビン（領域平均）。
方位は機体が\textbf{進む}向き。}
\begin{tabular}{lrrr}
\toprule
時期 & 平均距離 [\si{\kilo\metre}] & 方位 [度] & p90 [\si{\kilo\metre}] \\
\midrule
8 月 後半 & 50 & 281（西） & 85 \\
9 月 前半 & 51 & 32（北北東） & 83 \\
7 月 前半 & 54 & 96（東） & 89 \\
\midrule
2 月 前半 & 275 & 87（東） & 340 \\
1 月 後半 & 270 & 88（東） & 339 \\
12 月 後半 & 266 & 85（東） & 327 \\
\bottomrule
\end{tabular}
\end{table}

この表も和歌山とほとんど同じで、夏の飛行方位も同じように振れる。

\begin{itemize}
  \item 冬（12--2 月）は東へ \SIrange{248}{275}{\kilo\metre}。
  \item 7 月後半から 8 月は\textbf{西寄り}（方位 \ang{252}--\ang{281}）、
        9 月前半は北北東、7 月前半は東と、半月ごとに向きが変わる。
        夏は対流圏上部の弱い西風と成層圏の東風がほぼ釣り合うので、
        どちらが少し勝つかで向きが決まる。
  \item 7--8 月の平均距離は \SIrange{50}{58}{\kilo\metre} だが、
        東西に打ち消し合うので\textbf{平均落下点}は打ち上げ地点のすぐ近くになる。
        \SI{50}{\kilo\metre} 飛ぶ機体が、ある日は東へ、別の日は西へ落ちるということである。
\end{itemize}

最後の点は、平均落下点だけを見ると誤解しやすい。
たとえば 7 月前半の領域平均は、1 回あたり \SI{54}{\kilo\metre} 飛ぶのに、
平均落下点は東へ \SI{12}{\kilo\metre} しか動かない。
夏の計画では、平均の矢印ではなく散らばり（図\ref{fig:m8-ehime} の楕円）を見る必要がある。

\subsection{南レク御荘公園から上げた場合}

代表地点から上げたときの平均落下点を、陸か海か、最寄りの海岸まで何キロかで整理した。
海岸線は北海道の表\ref{tab:voising-offshore} と同じもの
（\texttt{frontend/public/geo/land.json}）を使っている。

\begin{table}[H]
\centering
\caption{南レク御荘公園から上げた場合の平均落下点。
飛行距離は 1 回ごとの距離の平均で、平均落下点までの距離とは違う。}
\label{tab:ainan-offshore}
\begin{tabular}{@{}llrrcr@{}}
\toprule
月 & ビン & {飛行距離 [\si{\kilo\metre}]} & {平均落下点の方位 [度]} & 陸／海 & {岸まで [\si{\kilo\metre}]} \\
\midrule
1 月  & 前半 & 264 & 88  & 海 & 59 \\
1 月  & 後半 & 270 & 88  & 海 & 54 \\
4 月  & 前半 & 177 & 88  & 海 & 32 \\
4 月  & 後半 & 155 & 87  & 海 & 24 \\
7 月  & 前半 & 53  & 95  & \textbf{陸} & 5 \\
7 月  & 後半 & 59  & 257 & 海 & 10 \\
10 月 & 前半 & 93  & 81  & 海 & 25 \\
10 月 & 後半 & 151 & 83  & 海 & 15 \\
\bottomrule
\end{tabular}
\end{table}

\textbf{7 月を除くと、平均落下点はどれも四国の南の海上である。}
東へ \SIrange{85}{265}{\kilo\metre} 流されて、土佐湾の沖から紀伊半島の南西沖に落ちる。
1 月は潮岬の西南西、岸から \SIrange{54}{59}{\kilo\metre} の沖で、
4 つの季節で最も岸から遠い。

\textbf{7 月は向きそのものが半月で入れ替わる。}
7 月前半の平均落下点は東へ \SI{9}{\kilo\metre} しか動かず、陸の上になる。
7 月後半は逆に西南西へ \SI{42}{\kilo\metre} 動き、豊後水道の南部の海上になる。
どちらも平均落下点が打ち上げ地点に近いのは、
上で述べたとおり東向きの日と西向きの日が打ち消し合った結果である。
1 回ごとの飛行距離は \SIrange{53}{59}{\kilo\metre} あり、
p90 では \SIrange{88}{96}{\kilo\metre} になる。

\begin{whybox}{北海道の VOISING と比べると}
VOISING は夏に岸から \SIrange{2}{4}{\kilo\metre}、他の季節は \SIrange{55}{80}{\kilo\metre} 沖に落ちる
（表\ref{tab:voising-offshore}）。
愛南町の冬の沖合いの距離（\SIrange{54}{59}{\kilo\metre}）は北海道と同じくらいだが、
春と秋は \SIrange{15}{32}{\kilo\metre} と北海道より岸に近い。
土佐湾の沖や室戸岬の南に落ちるためである。
夏は北海道が「岸のすぐそばの海」なのに対し、愛南町は
\textbf{陸にも海にも落ちうる}。回収の手配を両方に備える必要がある。

なお、この表は代表 1 地点の値である。
\texttt{out/csv/region\_sites\_ehime.csv} に「愛媛・高知」の 8 地点と
最寄り格子までの距離（\SIrange{12}{26}{\kilo\metre}）を出してあるので、
別の地点で見たい場合は \texttt{csv\_out.REFERENCE\_SITE} を変えて再実行すればよい。
\end{whybox}

\subsection{日変化と鉛直シアー}

\begin{figure}[H]
\centering
\includegraphics[width=0.75\linewidth]{ehime/F8_diurnal.pdf}
\caption[日変化]{夜（03 JST）と昼（15 JST）の平均風速の差を、季節ごとに高度に対して描いたもの
（\SI{300}{\hecto\pascal} 以下の層のみ、領域平均）。横軸の目盛りに注意してほしい。
差は 3 領域・4 季節のどれでも \SI{1}{\metre\per\second} 未満で、
日々の変動（数十 \si{\metre\per\second}）に比べて小さい。
これが\textbf{時刻を区別せず平均してよい根拠}である。
ビンを時刻で 4 分割すると標本数も 4 分の 1 になるので、
必要のないところで分割しない理由がここにある。}
\label{fig:f8-ehime}
\end{figure}

\begin{figure}[H]
\centering
\includegraphics[width=1.0\linewidth]{ehime/F9_shear.pdf}
\caption[鉛直シアー]{隣り合う気圧面のあいだの風速差を高度差で割ったもの。
機体そのものはシアーで壊れないが、大きいところでは
高度の推定が少しずれるだけで水平の予測が大きく動く。
最大はジェットの上下と地表付近に出る。}
\label{fig:f9-ehime}
\end{figure}


夜と昼の平均風速の差は最大 \SI{0.9}{\metre\per\second} で、
和歌山と同じく高度 \SI{1}{\kilo\metre} 付近（夜のほうが強い）に山がある。

鉛直シアーの最大は \SI{9.4}{(\metre\per\second)\per\kilo\metre}
（12 月後半、高度 \SI{16.3}{\kilo\metre}）で、
場所も大きさも和歌山（\SI{9.3}{(\metre\per\second)\per\kilo\metre}）とほぼ同じ。
冬の亜熱帯ジェットの真上で風が急に弱まるところである。

\subsection{既存シミュレータの ISA 仮定}

\begin{figure}[H]
\centering
\includegraphics[width=1.0\linewidth]{ehime/F10_isa_check.pdf}
\caption[ISA との比較]{気圧から高度への換算。左は JRA-3Q の実際の気圧面高度と、既存シミュレータが使う国際標準大気 ISA。右はその差。ISA は季節も緯度も天気も知らないので、季節によって系統的にずれる。}
\label{fig:f10-ehime}
\end{figure}


ISA からのずれは最大 \SI{790}{\metre}（夏）、冬は \SI{269}{\metre} である。
和歌山（夏 \SI{787}{\metre}）とほぼ同じで、北海道（夏 \SI{1036}{\metre}）より小さい。
